\PassOptionsToPackage{dvipsnames,svgnames,table}{xcolor}
\documentclass[11pt]{article}
\usepackage[a4paper,margin=22mm]{geometry}
\usepackage[no-math]{fontspec}
\setmainfont{DejaVu Serif}
\usepackage{amsmath,amssymb,amsthm,mathtools,graphicx,xcolor,hyperref,xurl,xspace,ifthen,enumerate}
\usepackage[shortlabels]{enumitem}
\setlength{\emergencystretch}{4em}
\AtBeginDocument{\special{pdf:mapline pzdr Dingbats <uzdr.pfb}}
\SetMathAlphabet{\mathbf}{normal}{OT1}{cmr}{bx}{n}
\SetMathAlphabet{\mathbf}{bold}{OT1}{cmr}{bx}{n}
\newtheorem{theo}{Theorem}[section]
\newtheorem{lemm}[theo]{Lemma}
\newtheorem{thm}[theo]{Theorem}
\newtheorem{proposition}[theo]{Proposition}
\newtheorem{corollary}[theo]{Corollary}
\newtheorem{definition}[theo]{Definition}
\newtheorem{theorem}[theo]{Theorem}
\newtheorem*{remas*}{Remarks}
\newtheorem{remas}[theo]{Remarks}
\newtheorem*{exam*}{Example}
\newtheorem{ques}[theo]{Question}
\providecommand{\og}{«\nobreakspace}
\providecommand{\fg}{\nobreakspace»}
\newtheorem{lem}[theo]{Lemma}
\newtheorem{lemma}[theo]{Lemma}
\newtheorem{prop}[theo]{Proposition}
\newtheorem{coro}[theo]{Corollary}
\newtheorem{corol}[theo]{Corollary}
\newtheorem*{conj*}{Conjecture}
\newtheorem{conj}[theo]{Conjecture}
\newtheorem{conjecture}[theo]{Conjecture}
\newtheorem{Proposition}[theo]{Proposition}
\newtheorem{theoreme}[theo]{Theorem}
\theoremstyle{definition}
\newtheorem{defi}[theo]{Definition}
\newtheorem{exam}[theo]{Example}
\newtheorem{exem}[theo]{Example}
\newtheorem{remark}[theo]{Remark}
\newtheorem{example}[theo]{Example}
\newtheorem{exams}[theo]{Examples}
\newtheorem{rema}[theo]{Remark}
\newtheorem{nota}[theo]{Notation}
\newtheorem{notation}[theo]{Notation}
\newtheorem*{notas}{Notations}
\newtheorem{result}[theo]{Result}
\newtheorem{question}[theo]{Question}
\newtheorem{prob}[theo]{Problem}
\newtheorem*{rema*}{Remark}
\newtheorem*{defi*}{Definition}
\newtheorem*{theo*}{Theorem}
\newtheorem*{lemm*}{Lemma}
\newtheorem*{prop*}{Proposition}
\newcommand{\enoncename}{}
\newtheorem{corpusenonce}[theo]{\enoncename}
\newenvironment{enonce}[2][]{\renewcommand{\enoncename}{#2}\begin{corpusenonce}}{\end{corpusenonce}}
\newenvironment{enonce*}[2][]{\par\medskip\noindent\textbf{#2.}\itshape}{\par\medskip}
\newenvironment{proofc}[1][\proofname]{\begin{proof}[#1]}{\end{proof}}
\newcommand{\Mc}{\mathcal{M}}
\newcommand{\pointir}{.\ ---\ }
\newcommand{\equalenv}[2]{\expandafter\let\csname#1\expandafter\endcsname\csname#2\endcsname\expandafter\let\csname end#1\expandafter\endcsname\csname end#2\endcsname}
\providecommand{\diagbox}[3][]{\begin{tikzpicture}[x=1em,y=1em,baseline=(current bounding box.center)]\useasboundingbox(0,0) rectangle(3,1.4);\draw(0,1.4)--(3,0);\node at(.5,.3){#2};\node at(2.5,1.1){#3};\end{tikzpicture}}
\providecommand{\eg}{e.g.\ }
\providecommand{\ie}{i.e.\ }
\providecommand{\cf}{cf.\ }
\providecommand{\longthanks}[1]{\par\smallskip\noindent #1\par\smallskip}
\providecommand{\qbinom}[2]{\genfrac{[}{]}{0pt}{}{#1}{#2}}

\providecommand{\mathof}[1]{\mathrm{#1}}

\numberwithin{equation}{section}


\def\endappendix{}


\usepackage{latexsym,amsmath,amssymb,verbatim,tikz,tikz-cd}
\usepackage{mathtools}
\usepackage{graphicx,xcolor}

\usepackage[english]{babel}
\usepackage{amssymb}
\usepackage{amsmath}
\usepackage{amsthm}
\usepackage{graphicx}
\usepackage{csquotes}
\usepackage{bm}
\usepackage{enumitem}
\usepackage{subcaption}
\usepackage{caption}
\renewcommand{\thesubfigure}{\Alph{subfigure}}
\usepackage{dirtytalk}
\providecommand{\nopunct}{\spacefactor1007\relax}
\newtheorem{problem}[theo]{Problem}

\newcommand{\sym}{\mathrm{Sym}}
\newcommand{\qsym}{\mathrm{QSym}}
\newcommand{\Des}{\operatorname{Des}}
\newcommand{\Asc}{\operatorname{Asc}}
\newcommand{\Syt}{\operatorname{SYT}}
\newcommand{\Short}{\operatorname{Short}}
\newcommand{\Loop}{\operatorname{Loop}}
\newcommand{\Stable}{\operatorname{Stable}}
\newcommand{\len}{\operatorname{len}}
\newcommand{\height}{\operatorname{height}}
\newcommand{\core}{\operatorname{core}}
\newcommand{\res}{\operatorname{res}}
\newcommand{\row}{\operatorname{row}}
\newcommand{\insrt}{\operatorname{insert}}
\newcommand{\UP}{\operatorname{UP}}
\newcommand{\DOWN}{\operatorname{DOWN}}


\newcommand{\ZZ}{\mathbb{Z}}
\newcommand{\NN}{\mathbb{N}}
\newcommand{\RR}{\mathbb{R}}
\newcommand{\QQ}{\mathbb{Q}}
\newcommand{\PP}{\mathbb{P}}

\newcommand{\F}{\mathcal{F}}
\newcommand{\Q}{\mathcal{Q}}
\newcommand{\A}{\mathcal{A}}
\newcommand{\M}{\mathcal{M}}
\newcommand{\I}{\mathcal{I}}
\newcommand{\PM}{\mathcal{PM}}
\newcommand{\Pcal}{\mathcal{P}}


\newcommand{\odds}{\mathrm{odds}}

\makeatletter
\newcommand{\proofstep}[1]{
  \par
  \addvspace{\medskipamount}
  \textbf{#1\@addpunct{.}}\enspace\ignorespaces
}

\newcommand\restr[2]{{
  \left.\kern-\nulldelimiterspace 
  #1 
  \vphantom{\big|} 
  \right|_{#2} 
  }}
\makeatother












\title{The set of matchings avoiding a short-chord-free pattern is closed under Knuth equivalence}
\author{Avichai Marmor}
\date{doi:10.5802/alco.351; CC BY 4.0}
\begin{document}
\maketitle
\noindent\textit{Editorial note.} The counted statement is the author's labelled lemma environment at publisher-source.tex lines 1300-1303; its proof environment opens at line 1304, immediately after \textbackslash{}end\{lemma\}, so the pairing is by adjacency and the retained proof range 1304-1307 is the whole environment.\par\medskip
\noindent\textit{Editorial note.} The definition def:pattern avoidance (source lines 1262-1279) is retained as context because the counted proof uses the notion of pattern containment introduced there; it is the only outside label of the retained proof.\par\medskip
\noindent\textit{Editorial note.} This article carries no label-number dump and the retained counter values are left empty for the whole source, so the printed numbering of the card is the wrapper's default; the identity of the counted item is its author label lem:avoiding set of short chord free matching is knuth like closed.\par\medskip
\noindent\textit{Editorial scope.} One original result of the pinned publisher article, counted once: closure under Knuth equivalence of the set of matchings avoiding a short-chord-free pattern (author label lem:avoiding set of short chord free matching is knuth like closed).\par\medskip
\noindent\textit{Allowed standard prerequisites (editorial).}\begin{itemize}
\item Pattern containment and avoidance. A matching m\_2 contains m\_1 when an order-preserving choice of vertices carries the chords and unmatched vertices of m\_1 to those of m\_2, and a set of matchings avoids m\_1 when none of its members contains it.. It is the notion appearing in the statement.
\item Elementary Knuth transformation. Knuth equivalence is generated by a local rewriting of three consecutive chords of a matching.. The proof only has to check one such transformation.
\item Short-chord-free pattern. The pattern m is assumed to have no short chord, so no chord of m joins two consecutive vertices.. It is what makes the case distinction in the proof exhaustive.
\end{itemize}\medskip
\begin{definition}
\label{def:pattern avoidance}
    Let $m_1 \in \M_{N_1}$ and $m_2 \in \M_{N_2}$ for some positive integers $N_1 \le N_2$. We say that $m_2$ \emph{contains the pattern} $m_1$, if there exist indices $1 \le i_1 < \dots < i_{N_1} \le N_2$ such that the following holds:
    \begin{itemize}
        \item For all $1 \le j < j' \le N_1$,
        \[
            (j,j') \in m_1 \Longleftrightarrow (i_j,i_{j'}) \in m_2.
        \]
        \item For all $1 \le j \le N_1$,
        \[
            (j) \in m_1 \Longleftrightarrow (i_j) \in m_2.
        \]
    \end{itemize}
    
    Otherwise, we say that $m_2$ \emph{avoids} $m_1$.

    For two sets of matchings $\M_1 \subseteq \M_{N_1}$ and $\M_2 \subseteq \M_{N_2}$, denote by $\M_2(\M_1)$ the set of matchings in $\M_2$ that avoid every matching in $\M_1$. In addition, denote $\M_2(m) := \M_2(\{m\})$ for a matching $m$.
\end{definition}

\begin{lemma}
\label{lem:avoiding set of short chord free matching is knuth like closed}
    Let $m \in \M_N$ be a short-chord-free matching, and let $N',f'$ be nonnegative integers. Then the pattern-avoiding set $\M_{N',f'}(m)$ is closed under Knuth equivalence.
\end{lemma}

\begin{proof}
    Let $m_1' \in \M_{N',f'}$ be a matching, and let $m_2' \in \M_{N',f'}$ be the result of applying an elementary Knuth transformation on $m_1'$. Assume that $m_1'$ contains the pattern $m$. It suffices to prove that $m_2$ contains $m$ too. By Definition~\ref{def:pattern avoidance}, there exist indices $1 \le i_1 < \dots < i_{N} \le N'$ such that $(j,j') \in m \Longleftrightarrow (i_j,i_{j'}) \in m_1'$ for all $1 \le j < j' \le N$ and $(j) \in m \Longleftrightarrow (i_j) \in m_1'$ for all $1 \le j \le N$. Without loss of generality, we may assume that $(i,i+1) \in m_1'$ and $m_2'$ is obtained from $m_1'$ by interchanging the chord $(i,i+1)$ with the vertex $i+2$. The pattern $m$ is short-chord-free, so $i, i+1 \notin \{i_1,\dots,i_N\}$. Therefore,
    by considering the index set $\{i_1,\dots,i_N\}$ if $i+2 \notin \{i_1,\dots,i_N\}$ and the set $\{i\}\cup\{i_1,\dots,i_N\} \setminus \{i+2\}$ otherwise, we obtain that the matching $m_2'$ contains $m$ as well.
\end{proof}

\end{document}
